\documentclass[10pt,letterpaper]{article}
\usepackage{cornell-notes}

\title{Clause 17.06.04.05: Get New Handler}
\author{}
\date{}
\setDocTitle{Clause 17.06.04.05: Get New Handler}
\setDocSubtitle{C++ 2024}
\setDocOwner{C++ 2024 Cornell Notes}
\setDocDate{\today}
\setCornellCollection{C++ 2024 Cornell Notes}
\setCornellUnitType{Clause}
\setCornellUnitNumber{17.06.04.05}
\setCornellUnitTitle{Get New Handler}
\hypersetup{pdftitle={Clause 17.06.04.05: Get New Handler},pdfsubject={Cornell notes},pdfauthor={}}

\begin{document}
\maketitle
\begin{CornellOverviewBox}{Overview}
\textbf{Classification:} Standard library - Normative\\[2pt]
\textbf{Learning objective:} Understand the rules, terminology, and semantic consequences associated with get\_new\_handler.
\end{CornellOverviewBox}\begin{CornellSummaryBox}{Summary}
{\sffamily\bfseries\color{Accent} Summary}\par\vspace{3pt}Section 17.6.4.5 focuses on get\_new\_handler. Apply the specified interface together with its mandates, effects, result conditions, exception behavior, and complexity constraints. When applying it, preserve the distinction between mandatory requirements, recommendations, permissions, and behavior deliberately left undefined, unspecified, or implementation-defined.
\end{CornellSummaryBox}

\end{document}
